\chapter{SIMPULAN DAN SARAN}

\section{Simpulan}

% TODO: Tuliskan simpulan dalam bentuk paragraf yang efektif. Simpulan harus:
% (1) menjawab tujuan dan rumusan masalah pada Bab 1;
% (2) menjelaskan permasalahan yang ingin dipecahkan, metode yang dipakai,
%     dan alasan penggunaan metode tersebut;
% (3) menjelaskan hasil implementasi secara singkat, padat, dan jelas;
% (4) menunjukkan kekurangan penelitian yang dapat disarankan untuk
%     penelitian lanjutan.

Berdasarkan penelitian yang telah dilakukan, dapat diambil simpulan sebagai berikut. Lorem ipsum dolor sit amet, consectetur adipiscing elit. Sed do eiusmod tempor incididunt ut labore et dolore magna aliqua. Ut enim ad minim veniam, quis nostrud exercitation ullamco laboris nisi ut aliquip ex ea commodo consequat. Duis aute irure dolor in reprehenderit in voluptate velit esse cillum dolore eu fugiat nulla pariatur.

\section{Saran}

% TODO: Tuliskan saran untuk penelitian lanjutan berdasarkan kekurangan atau
% keterbatasan penelitian yang dilakukan.

Berdasarkan hasil penelitian, saran untuk pengembangan selanjutnya adalah:

\begin{enumerate}
    \item Lorem ipsum dolor sit amet, consectetur adipiscing elit.
    \item Sed do eiusmod tempor incididunt ut labore et dolore magna aliqua.
    \item Ut enim ad minim veniam, quis nostrud exercitation ullamco laboris.
\end{enumerate}
